\documentclass[9pt]{elife}
\usepackage{amsthm}
\usepackage{amsmath,amssymb,mathtools}
\usepackage{siunitx}
\usepackage{graphicx}
\usepackage{booktabs}
\usepackage{hyperref}
\usepackage{tcolorbox}
\usepackage{bm}

\title{MacroIR: A Bayesian framework for analytically exact interval likelihoods in time-averaged biological data}

\author[1*]{Luciano Moffatt}
\affil[1]{Instituto de Qu\'imica F\'isica de los Materiales, Medio Ambiente y Energ\'ia (INQUIMAE), CONICET; Facultad de Ciencias Exactas y Naturales, Universidad de Buenos Aires, Ciudad de Buenos Aires, C1428EHA, Argentina}
\corr{lmoffatt@qi.fcen.uba.ar}

\begin{document}
	\maketitle
	\begin{abstract}
		Time-averaged observations are pervasive in experimental biology, yet most inference frameworks assume instantaneous sampling and thus misestimate systems whose dynamics evolve faster than the acquisition window. 
		\textbf{MacroIR} (Macroscopic Interval Recursive) formulates a fully Bayesian likelihood for interval-averaged signals generated by continuous-time Markov processes. 
		Exploiting the Markov property, the method treats each observation window as a bridge between an initial and a final state. 
		Given the prior at the start and the kinetic parameters, this pair suffices to predict the probability of all trajectories within the interval. 
		The observed average depends jointly on both states, enabling a Bayesian update over the Cartesian distribution of start–end pairs using the analytical mean and variance of the observable. 
		Marginalizing over the initial state yields the posterior for the next interval, completing a deterministic recursion without Monte Carlo sampling. 
		In the linear–Gaussian limit, MacroIR reduces exactly to the Kalman filter, clarifying the latter as a special case of Bayesian inference on Markov processes. 
		This formulation unifies stochastic kinetics, time-averaged observables, and Bayesian evidence evaluation, providing unbiased parameter and state inference across molecular and cellular datasets.
	\end{abstract}
	\begin{abstract}
		Time-averaged measurements dominate experimental biology, yet most inference methods assume instantaneous observations, yielding biased estimates when intrinsic dynamics outpace acquisition. 
		\textbf{MacroIR} (Macroscopic Interval Recursive) provides a fully Bayesian formulation of the likelihood for interval-averaged signals generated by continuous-time Markov processes. 
		By conditioning on both the start and end states of each sampling window, MacroIR computes the interval mean and variance in closed form and recursively updates the posterior probability distribution over hidden states. 
		Under a Gaussian-moment approximation, these updates reduce to the familiar Kalman recursion, but the Bayesian derivation remains valid beyond the linear–Gaussian regime. 
		The framework unifies ensemble kinetics, time-averaged observables, and Bayesian evidence computation. 
		Applications to simulated and experimental data demonstrate that MacroIR eliminates structural bias in time-averaged biological recordings and enables reliable inference of underlying dynamics across molecular, cellular, and behavioral timescales.
	\end{abstract}
	
	\section*{Introduction}
	
	Time-averaged measurements are ubiquitous in experimental biology—electrophysiological amplifiers integrate currents over microseconds, fluorescence detectors average photon flux over milliseconds, and population assays aggregate signals over hours \citep{Hille2001, Denk1990, Elowitz2002}. 
	Yet most inference methods treat such data as instantaneous samples. 
	When intrinsic dynamics evolve on comparable timescales, this mismatch between physical averaging and model assumptions distorts likelihoods and biases inferred parameters \citep{ColquhounHawkes1982, Qin2000, Munsky2009, Moffatt2007}.
	
	Continuous-time Markov models (CTMCs) remain the standard framework for describing stochastic molecular and cellular dynamics. 
	Over recent decades, inference methods have evolved from rate-equation and hidden-Markov formulations to recursive Bayesian estimators that continuously update knowledge of the system’s state \citep{Csanady2006, Siekmann2011, Bronson2013}.  
	Münch \emph{et al.} (2022) recently introduced a generalized Bayesian filter for ion-channel kinetics that unifies current and fluorescence data and explicitly integrates mean dynamics over finite sampling windows \citep{Munch2022}.  
	However, their approach retains an implicit Gaussian approximation for the posterior distribution and assumes instantaneous observation within each bin, leading to bias when acquisition windows approach the system’s dwell times.
	
	A direct precursor of this work was the \emph{macroscopic recursive} (Macro R) algorithm introduced by Moffatt (2007) for the Bayesian estimation of kinetic rates from macroscopic current fluctuations.  
	Macro R formulated the full recursive likelihood for ensemble recordings, approximating the hidden state distribution by a Gaussian and updating its moments after each measurement.  
	The resulting recursion was mathematically identical to a Kalman filter but conceptually derived from first Bayesian principles.  
	The present framework, \textbf{MacroIR} (Macroscopic Interval Recursive), extends that Bayesian theory to \textbf{time-averaged observations}, deriving analytically exact interval likelihoods and recursive updates for the posterior distribution.  
	This extension restores consistency between the physical measurement process and the probabilistic model, allowing unbiased inference and rigorous model comparison via Bayesian evidence.
	
	\section*{Theory}
	
	\subsection*{From the macroscopic recursive (Macro R) to the interval-recursive (Macro IR) formulation}
	
	In the Macro R approach \citep{Moffatt2007}, the ensemble of channel occupancies is represented by a probability distribution $p(\mathbf{s},t)$ over discrete states that evolves under a rate matrix $Q$.  
	After each instantaneous measurement $y_t$, Bayes’ rule updates the posterior:
	\[
	p^{\text{post}}(\mathbf{s}|y_t) \propto p(y_t|\mathbf{s})\,p^{\text{prior}}(\mathbf{s}).
	\]
	When the ensemble contains many channels, $p(\mathbf{s})$ can be approximated by a multivariate normal characterized by mean vector $\bm{\mu}$ and covariance matrix $\Sigma$.  
	The propagation of these moments between observations yields what is commonly known as the Kalman recursion but here interpreted as the \emph{Gaussian-moment limit} of Bayesian inference.
	
	MacroIR generalizes this logic to time-averaged observations.  
	Instead of observing $y_t$ at a single point, detectors measure the mean observable over an acquisition window $[0,t]$:
	\[
	\bar y_{0\to t} = \frac{1}{t}\int_0^t \gamma(x_s)\,ds.
	\]
	The likelihood must therefore condition on both the start and end states of the hidden Markov process within each interval.
	
	\subsection*{Boundary-conditioned Bayesian formulation (meta-state representation)}
	
	Let $x_t$ denote a CTMC with generator $Q$ and propagator $P(t)=e^{Qt}$.  
	For a linear observable $\bm{\gamma}\in\mathbb{R}^k$, the conditional mean of the interval average, given start state $i$ and end state $j$, is
	\begin{equation}
		\gamma_{i\to j}(t)
		= \frac{1}{P_{i\to j}(t)}
		\sum_{n_1,n_2}
		V_{i n_1}\,(V^{-1}\Gamma V)_{n_1 n_2}\,V^{-1}_{n_2 j}\,
		E_2(\lambda_{n_1} t, \lambda_{n_2} t),
	\end{equation}
	where $Q=V\Lambda V^{-1}$, $\Gamma=\mathrm{diag}(\bm{\gamma})$, and
	\[
	E_2(x,y)=\frac{e^x-e^y}{x-y}.
	\]
	An analogous expression for the variance involves a three-exponential kernel $E_3(x,y,z)$ derived from the spectral decomposition of $Q$:
	\begin{align}
		\mathrm{Var}(\Gamma_{i\to j}) 
		= E[\Gamma^2_{i\to j}] - (E[\Gamma_{i\to j}])^2,
	\end{align}
	\begin{align}
		E[\Gamma^2_{i\to j}] 
		= \frac{1}{P_{i\to j}(t)}
		\sum_{k_1,k_2,n_1,n_2,n_3}
		V_{i n_1}V^{-1}_{n_1 k_1}\gamma_{k_1}
		V_{k_1 n_2}V^{-1}_{n_2 k_2}\gamma_{k_2}
		V_{k_2 n_3}V^{-1}_{n_3 j}
		E_3(\lambda_{n_1}t,\lambda_{n_2}t,\lambda_{n_3}t).
	\end{align}
	
	\begin{tcolorbox}[colback=gray!5,colframe=gray!60,title=Boundary-conditioned two-point states (“meta-states”)]
		Define the ordered pair $(i,j)$ representing start and end states over duration $t$.  
		The $k^2$ possible pairs form a space whose generator is the Kronecker sum
		$Q\oplus Q = Q\otimes I + I\otimes Q$.  
		Because the propagator factorizes as $e^{(Q\oplus Q)t}=e^{Qt}\otimes e^{Qt}$, the integrals for the interval mean and variance reduce to bilinear forms in the original $k$-dimensional state space—no explicit $k^2$ object is needed.
	\end{tcolorbox}
	
	This boundary-conditioned Bayesian formulation provides the exact likelihood of observing the interval-averaged measurement $\bar y_{0\to t}$ given $(i,j)$.  
	In the limit $t\!\to\!0$, $\gamma_{i\to j}(t)\!\to\!\gamma_i$ and the formulation collapses to the instantaneous Macro R likelihood.
	
	\subsection*{Gaussian-moment Bayesian recursion}
	
	For an ensemble of $N_{\mathrm{ch}}$ independent channels, the joint occupancy distribution can be approximated by a multivariate normal with mean $\bm{\mu}$ and covariance $\Sigma$.  
	The first two moments of the meta-state Bayesian recursion yield a tractable Gaussian-moment update for the posterior distribution.
	
	Given a prior $(\bm{\mu}^{\mathrm{prior}}_0,\Sigma^{\mathrm{prior}}_0)$ and an observed interval-averaged current $y^{\mathrm{obs}}_{0\to t}$, the predicted mean and variance of the observable are
	\begin{align}
		y^{\mathrm{pred}}_{0\to t} &= N_{\mathrm{ch}}\,(\bm{\mu}^{\mathrm{prior}}_0\!\cdot\!\bm{\gamma}_0),\\
		\sigma^2_{y^{\mathrm{pred}}_{0\to t}} &= 
		\frac{\epsilon^2}{t} + \nu^2_{\mathrm{white}} + \nu^2_{\mathrm{pink}}
		+ N_{\mathrm{ch}}\!\left[\bm{\gamma}_0^\top(\Sigma^{\mathrm{prior}}_0-\mathrm{diag}\,\bm{\mu}^{\mathrm{prior}}_0)\bm{\gamma}_0 
		+ \bm{\mu}^{\mathrm{prior}}_0\!\cdot\!\mathbf{E}\right],
	\end{align}
	where $\bm{\gamma}_0=\{\gamma_{i\to j}(t)\}$ are the interval-averaged conductances derived above.  
	
	The posterior moments follow directly from Bayes’ rule under Gaussian closure:
	\begin{align}
		\bm{\mu}^{\mathrm{post}}_{t} &=
		\bm{\mu}^{\mathrm{prior}}_0 P(t)
		+ \frac{y^{\mathrm{obs}}_{0\to t}-y^{\mathrm{pred}}_{0\to t}}{\sigma^2_{y^{\mathrm{pred}}_{0\to t}}}\,K^\mu_{0\to t},\\
		\Sigma^{\mathrm{post}}_{t} &=
		\Sigma^{\mathrm{prior}}_{t}
		- \frac{1}{\sigma^2_{y^{\mathrm{pred}}_{0\to t}}}
		(K^\mu_{0\to t})^\top K^\mu_{0\to t},
	\end{align}
	with the optimal linear estimator
	\[
	K^\mu_{0\to t} = \frac{\mathrm{Cov}(x_t,\bar y_{0\to t})}{\mathrm{Var}(\bar y_{0\to t})}.
	\]
	When the system is linear and noise is Gaussian, these equations coincide with the classical Kalman filter; otherwise, they remain valid as an approximate Bayesian recursion.
	
	\section*{Methods}
	
	\subsection*{Synthetic models and data}
	We simulate $k$-state CTMCs with known $Q$ and linear observables $\bm{\gamma}$. 
	Traces are averaged over bins $\Delta$ spanning sub- to supra-kinetic regimes. 
	White and pink noise components $(\epsilon^2,\nu^2)$ follow baseline experimental estimates. 
	Priors on rate constants are log-uniform, and noise parameters use weakly-informative priors.
	
	\subsection*{Comparative benchmarks}
	We implemented three alternatives: (i) midpoint likelihood (single-point per bin), (ii) non-recursive interval likelihood (NR) without boundary conditioning, and (iii) instantaneous-sample approximation. 
	Each method was compared in terms of bias, RMSE, coverage of 95\% credible intervals, and runtime. 
	MacroIR used identical priors and samplers. 
	Model ranking employed Bayesian evidence computed via thermodynamic integration.
	
	\subsection*{Score test}
	Let $S(\theta)=\nabla_\theta \log \mathcal{L}(\theta)$ denote the per-interval score. 
	Under a correct likelihood, $\mathbb{E}[S]=0$ and $\mathrm{Var}(S)=\mathcal{I}(\theta)$ (Fisher information). 
	We estimated $\langle S\rangle$ and $\mathrm{Var}(S)$ across intervals and compared them to plug-in Fisher estimates from MacroIR, evaluating deviations as a measure of informational self-consistency.
	
	\section*{Results}
	
	\subsection*{Synthetic recovery and benchmark comparisons}
	MacroIR recovered unbiased parameters across all $\Delta/\tau_{\min}$ regimes, with calibrated uncertainty and higher evidence for true models. 
	Midpoint and NR approaches displayed increasing bias and under-coverage as $\Delta$ grew, while instantaneous approximations failed when $\Delta\gtrsim\tau_{\min}$.
	
	\subsection*{Score-based self-consistency}
	Only MacroIR yielded mean score $\approx0$ with variance matching Fisher predictions across sampling regimes. 
	Alternative methods showed nonzero mean scores and underestimated information, especially at large $\Delta$, revealing internal inconsistency between their assumed likelihood and the physical averaging process.
	
	\subsection*{Biological case study: P2X receptor kinetics}
	Reanalysis of ultrafast ATP-pulse macroscopic currents from P2X$_2$ receptors showed that MacroIR reproduced experimental responses and correctly ranked candidate kinetic schemes, revealing asymmetric coupling between ligand binding and subunit rotation. 
	These data were previously analyzed using the Macro R approximation \citep{Moffatt2007}; MacroIR reproduces those results in the instantaneous limit and reveals additional asymmetry once interval averaging is incorporated.
	
	\subsection*{Computational performance}
	Computation scales linearly with the number of intervals and as $O(k^3)$ with state dimension, dominated by matrix exponentials and linear solves. 
	Including an explicit analog filter via state-space augmentation increases runtime $\sim$10$\times$ but remains tractable for moderate $k$.
	
	\section*{Discussion}
	
	MacroIR provides a unified Bayesian framework for inference from time-averaged observations governed by continuous-time Markov processes. 
	It generalizes the macroscopic recursive algorithm of Moffatt (2007), which already cast macroscopic current analysis in Bayesian terms, to handle explicit temporal averaging and to compute model evidence. 
	The resulting equations encompass the Kalman recursion as a limiting case but retain their validity beyond linear–Gaussian assumptions.  
	By restoring consistency between the physics of sampling and the statistical model, MacroIR eliminates the structural bias intrinsic to instantaneous-likelihood approaches.
	
	The Bayesian foundation of MacroIR allows immediate extensions to non-Gaussian noise, hierarchical priors, and mixture models, and provides a natural route to Bayesian evidence comparison across kinetic schemes. 
	This generality distinguishes it from purely algorithmic filters and repositions the theory of macroscopic inference within the domain of full probabilistic reasoning.
	
	\subsection*{Limitations and outlook}
	(1) Scaling may become demanding for very large $k$; sparsity or exponential-Krylov methods could improve performance. 
	(2) Explicit analog filtering requires state-space augmentation or kernel convolution and increases runtime. 
	(3) Robustness to model misspecification and application to hierarchical or mixture CTMCs remain open avenues. 
	(4) Nonlinear observables will require moment-closure or variational approximations beyond the Gaussian level.
	
	\section*{Data and Code Availability}
	Code and scripts to reproduce all results will be released publicly (GitHub/Zenodo) at submission, including minimal working examples and figure-generation notebooks.
	
	\section*{Acknowledgements}
	I thank colleagues for discussions on Bayesian conditioning of time-averaged processes and the INQUIMAE computing facility for support.
	
	\section*{Author Contributions}
	LM conceived the study, developed the theory, implemented algorithms, analyzed data, and wrote the manuscript.
	
	\bibliography{biblio}
	
\end{document}
